\documentclass[12pt,twoside]{report}
\tolerance=1000
\hbadness=10000
\raggedbottom
%\usepackage[colorlinks,citecolor=blue,urlcolor=blue,bookmarks=false,hypertexnames=true]{hyperref}
%\RequirePackage{doi}

\usepackage[printonlyused]{acronym}
\usepackage{hyperref}
%\usepackage{doi}
%\usepackage{fancyref}
%\usepackage{nameref}
\usepackage[flushleft]{threeparttable}
\usepackage{rotating}
\usepackage{graphicx}
\usepackage{graphics}
\usepackage{url}
\usepackage{subfigure}
\usepackage{amsfonts}
\usepackage{amsmath,latexsym,amssymb,amsthm}
\usepackage[mathscr]{eucal}
\usepackage{verbatim}
\usepackage{multirow}
\usepackage{multicol}
\usepackage{setspace}
\usepackage{lscape,fancyhdr}
\usepackage{longtable}
\usepackage{lipsum}
\usepackage{titletoc}
\usepackage{varwidth}
\usepackage[utf8]{inputenc}
\usepackage{longtable}
%\usepackage[square, numbers, comma, sort&compress]{natbib}
%\usepackage{algorithm, algpseudocode}
%\usepackage[toc,nonumberlist]{glossaries}
\usepackage[nottoc]{tocbibind}  %to appear bibliography in TOC
\usepackage{a4}
\usepackage[top=4cm, left=4cm, right=2.5cm]{geometry}
%%%%%%%%%%%%%%%%%%%%%%%%%%%%%%%%% EXTRA %%%%%%%%%%%%%%%%%%%%%%%%
\usepackage{xparse}
\usepackage{mathtools}
\usepackage{romannum}
\usepackage{algorithm2e}
\usepackage{algpseudocode}
\usepackage{amsfonts}
\usepackage{wrapfig}
\usepackage{bm}
\usepackage{emptypage}%This will automatically suppress page numbers on all inserted blank pages (like those after \cleardoublepage).
\usepackage{verbatim}
\usepackage{multirow}
\usepackage{ragged2e}
\usepackage{enumerate}
\usepackage{enumitem}
\usepackage{tikz}
\usepackage{tikzpeople}
\usepackage{tkz-berge}
\usetikzlibrary{shapes,snakes,positioning,shapes.misc,arrows.meta}
\usepackage[super]{nth}
\usepackage{float}
\usepackage{url}
\usepackage{cite}
%\usepackage{booktabs} %shortcuts for pre-defined table rule commands
%\glsdisablehyper   % disable hyperlink in glossary
\DeclareUnicodeCharacter{2212}{-}
%\DeclareUnicodeCharacter{00A0}{}
%%%%%%%%%%%%%%%%%%%%%%%%%%%%%%%%%%% KP %%%%%%%%%%%%%%%%%%%%%%%
\usepackage{pgfplots,pgfplotstable}
\pgfplotsset{width=10cm,compat=1.9}
\usetikzlibrary{shadows,calc}
\usetikzlibrary{calc,shapes.multipart,chains,arrows}
\usetikzlibrary{matrix}
\usetikzlibrary{automata,positioning}
\usetikzlibrary{positioning}
\usetikzlibrary{decorations.markings,decorations.pathmorphing,calc, intersections}
\pgfplotsset{}
\renewcommand{\arraystretch}{1.25}
%%%%%%%%%%%%%%%%%%%%%%%%%%%%%%%%%%%%%%%%%%%%%%%%%%%%%%%%%%%%%%%
%\setlength{\textwidth}{14.7cm}
%\renewcommand{\baselinestretch}{1.05}
% Set up proclamation styles
\newtheorem{secn}{Definition}[section]
\newtheorem{theorem}[secn]{Theorem}
\newtheorem{corollary}[secn]{Corollary}
\newtheorem{proposition}[secn]{Proposition}
\newtheorem{lemma}[secn]{Lemma}
\newtheorem{definition}[secn]{Definition}
\newtheorem{example}[secn]{Example}
\newtheorem{remark}[secn]{Remark}
\newtheorem{problem}[secn]{Problem}
%\newtheorem{algorithm}[secn]{Algorithm}
\newtheorem{construction}[secn]{Construction}
\newcommand{\pf}{{\bf Proof:}}
\newcommand{\ai}{{\sf AI}}

\setcounter{section}{0}

\newcommand\floor[1]{\lfloor#1\rfloor}
\newcommand\ceil[1]{\lceil#1\rceil}

%Numbering within environments
\newcommand{\newsection}[1]{\setcounter{equation}{0} \setcounter{secn}{0}\section{#1}}
\numberwithin{equation}{section}
\renewcommand{\theequation}{\thesection.\arabic{equation}}
\renewcommand{\qed}{\hfill \rule{2mm}{2mm}}

%%%%%%%%%%%%%%%%%%%%%%%%%%%%%%%%%%%%%%%%%%%%%%%%%%%%%%%%%%%
%%%%%%%%%%%%%%%%%%%%%%%%%%%%%%%%%%%%%%%%%%%%%%%%%%%%%%%%%%%%%%%
\newcommand{\hll}[1]{\colorbox{yellow}{$\displaystyle #1$}}
%\newtheorem{theorem}{Theorem}
%\newtheorem{proof}{Proof}
%\newtheorem{lemma}{Lemma}
\newtheorem{mydef}{Definition}
\newtheorem{prop}{Proposition}
\newtheorem{prf}{Proof}
%%%%%%%%%%%%%%%%%%%%%%%%%%%%%%%%%%%%%%%%%%%%%%%%%%%%%%%%%%%%%%%

%%%%%%%%%%%%%%%%%%%%%%%%%%%%%%%%%%%%%%%%%%%%%%%%%%%%%%%%%%%
\usetikzlibrary{shapes, arrows, calc, arrows.meta, fit, positioning} % these are the parameters passed to the library to create the node graphs
\tikzset{-Latex,auto,node distance =3.8 cm and 3.8 cm, thick,% node distance is the distance between one node to other, where 1.5cm is the length of the edge between the nodes
state/.style ={ellipse, draw, minimum width = 0.9 cm}, % the minimum width is the width of the ellipse, which is the size of the shape of vertex in the node graph
point/.style = {circle, draw, inner sep=0.18cm, fill, node contents={}},
bidirected/.style={Latex-Latex,dashed}, % it is the edge having two directions
connected/.style={dashed},
el/.style = {inner sep=2.5pt, align=right, sloped}
}
%%%%%%%%%%%%%%%%%%%%%%%%%%%%%%%%%%%%%%%%%%%%%%%%%%%%%%%%%%%
\begin{document}
    \thispagestyle{empty}
%\cleardoublepage
\pagestyle{empty}
\vspace*{-1.75cm}
\begin{center}
\begin{Large}
{\textbf{Title of the B.Tech./Executive M.Tech. Thesis}}\\
\end{Large}
\end{center}
\vspace{0.2cm}
\begin{center}
 \begin{large}
 Thesis  submitted to \\
 \vspace{.2cm}
 \large{\textbf{Indian Institute of Information Technology Kalyani}}\\
 \vspace{.1cm}
 in partial fulfilment of the requirements\\
 \vspace{0.1cm}
 for the award of the degree\\
 \vspace{0.3cm}
 of\\
 \vspace{0.3cm}
 \textbf{Bachelor of Technology/Executive Master of Technology}
 \vspace{0.3cm}
 \\in\\
 \vspace{0.3cm}
 \textbf{Computer Science and Engineering/Artificial Intelligence and Data Science}
 \end{large}
\end{center}

\vspace{0.00cm}
 \begin{center}
 \begin{large}
 by\\
 \vspace{0.2cm}
 \textbf{Name of the Student}\\
 \vspace{0.1cm}
 Registration No.: \\
 \end{large}
\end{center}

\vspace{0.1cm}

\begin{center}
 \begin{tabular}{l}
 \includegraphics[width = 3.5 cm]{Figure/header/IIITK.png}
 \end{tabular}
\end{center}
\vspace{-0.3cm}
\begin{center}
\begin{large}
 \vspace{0.25cm}
 \hspace{-15mm} DEPARTMENT OF COMPUTER SCIENCE AND ENGINEERING\\
 \hspace{-15mm} INDIAN INSTITUTE OF INFORMATION TECHNOLOGY KALYANI\\
 KALYANI - 741235, WEST BENGAL, INDIA\\
 \vskip 50pt
 May, 2025
\end{large}
\end{center}


\newpage
\thispagestyle{empty}
\cleardoublepage
\pagestyle{empty}
\vspace*{-1.75cm}
\begin{center}
\begin{Large}
{\textbf{Title of the B.Tech./Executive M.Tech. Thesis}}\\
\end{Large}
\end{center}
\vspace{0.2cm}
\begin{center}
 \begin{large}
 Thesis  submitted to \\
 \vspace{0.2cm}
 \large{\textbf{Indian Institute of Information Technology Kalyani}}\\
 \vspace{0.1cm}
 in partial fulfilment of the requirements\\
 \vspace{0.1cm}
 for the award of the degree\\
 \vspace{0.3cm}
 of\\
 \vspace{0.3cm}
 \textbf{Bachelor of Technology/Executive Master of Technology}
 \vspace{0.3cm}
 \\in\\
 \vspace{0.3cm}
 \textbf{Computer Science and Engineering/Artificial Intelligence and Data Science}
 \end{large}
\end{center}

\vspace{-0.2cm}
\begin{center}
 \begin{large}
 by\\
 \vspace{0.2cm}
 \textbf{Name of the student}\\
 \vspace{0.1cm}
 Registration No.: PhD/R/CSE/19008/510\\
 \end{large}
\end{center}

\vspace{0.1cm}
\begin{center}
 \begin{large}
 under the supervision of\\
 \vspace{.2cm}
 \textbf {Name of the Supervisor}\\
 \vspace{.2cm}
 \end{large}
\end{center}
\begin{center}
 \begin{tabular}{l}
 \includegraphics[width = 3.5 cm]{Figure/header/IIITK.png}
 \end{tabular}
\end{center}
\vspace{-0.3cm}
\begin{center}
\begin{large}
  \vspace{0.15cm}
  \hspace{-15mm} DEPARTMENT OF COMPUTER SCIENCE AND ENGINEERING\\
  \hspace{-15mm} INDIAN INSTITUTE OF INFORMATION TECHNOLOGY KALYANI\\
   KALYANI - 741235, WEST BENGAL, INDIA\\
   \vskip 50pt
   May, 2025
\end{large}
\end{center}

%\setcounter{secnumdepth}{3}
%\setcounter{tocdepth}{2}
\newpage
\cleardoublepage
\chapter*{}
\thispagestyle{empty}
\vspace{6cm}
\begin{center}
{\Large\Large\Large \textit{Dedicated to my ......... It is
optional}}
\end{center}

\newpage
 \thispagestyle{empty} \mbox{}
 \pagenumbering{roman}
 \setcounter{page}{0}
 \pagestyle{plain}

%\input{certi_approv}
\cleardoublepage
%%\input{cert1.tex} [take a print of the text written in template on Institute letterhead, include scanned copy as shown in cerrt1.tex, in place of ch-certificate_sup]
  \cleardoublepage
 \thispagestyle{empty}
 \begin{center}
 {\bf\uppercase{Indian Institute of Information Technology Kalyani}}\\
 \vspace{-0.3cm}
 \end{center}
%\begin{center}
\begin{multicols}{2}
%\vspace{-1cm}
 \flushleft
 \vspace{-2cm}
 \includegraphics[width=15mm]{Figure/header/IIITK.png}
 \columnbreak
 \par
 \begin{center}
 \hspace{-75mm} {\large An Institute of National Importance \vspace{0.15cm}}\\
 \hspace{-73mm} \textmd{(Autonomous Institution under MoE, Govt. of India \& \vspace{0.15cm}} \\
 \hspace{-66mm}\textmd{Department of Information Technology \& Electronics, Govt. of West Bengal)}\\
 \hspace{-73mm}\textmd{WEBEL IT Park, Opposite of Kalyani Water Treatment Plant \vspace{0.1cm}}\\
 \hspace{-71mm} \textmd{Near Buddha Park, Dist. Nadia, P.O. Kalyani - 741235, West Bengal \vspace{0.1cm}}\\
 \hspace{-73mm} \textmd{Email: office@ilitkalyani.ac.in, Website: https://iiitkalyani.ac.in/}
 \end{center}
\end{multicols}
 \vspace{-1cm}
 \begin{center}
\rule{\textwidth}{0.04cm}
\end{center}
 \vspace{0.5cm}
\begin{center}
{\LARGE \textsc{Certificate}}
\end{center}
\vspace{0.5cm} This is to certify that the thesis entitled
``\textbf{Title of the Thesis}'', being submitted by \textbf{Name of
the Student} (\textbf{Reg. No:}), a undergraduate/postgraduate
student at the \textbf{Indian Institute of Information Technology
Kalyani}, West Bengal, India, for the award of the \textbf{Bachelor
of Technology}/\textbf{Executive Master of Technology} in
\textbf{Computer Science and Engineering}/\textbf{Artificial
Intelligence and Data Science}, is an original research work carried
out by him under my supervision and guidance.

\noindent \par The thesis has fulfilled all the requirements as per
the regulations of \textbf{IIIT Kalyani} and, in my opinion, has
attained the standards required for submission. The work,
techniques, and results presented herein have not been submitted to
any other university or institute for the award of any other degree
or diploma.

\vspace{1in}

\begin{flushright}
%    \includegraphics[scale=0.4]{fig.jpg}\\
%    \vskip-3mm
    \textbf{(Name of the Guide)}\\
    \text{Assistant Professor}\\
    \text{Department of Computer Science and Engineering}\\
    \text{Indian Institute of Information Technology Kalyani}\\
    \text{West Bengal - 741235, India}
\end{flushright}

 \cleardoublepage
  \cleardoublepage
 \thispagestyle{empty}
 \begin{center}
  {\LARGE \textsc{Declaration}}
 \end{center}
 \vspace{0.3cm}

%\begin{flushright}
%\begin{tabular}{l}
%"\textit{The time has come,}" \textit{the Walrus said,}\\
%"\textit{To talk of many things}" \dots \\
%\hspace{1in} - {\scriptsize Lewis Carroll, Through the Looking-Glass}\\
%\end{tabular}
%\end{flushright}

I hereby certify that the work which is being presented in the
thesis entitled ``\textbf{Title of the Thesis}'' in the partial
fulfillment of the requirements for the award of the degree of
\textbf{Bachelor of Technology}/\textbf{Executive Master of
Technology} in \textbf{Computer Science and
Engineering}/\textbf{Artificial Intelligence and Data Science} in
the \textbf{Department of Computer Science and Engineering},
\textbf{Indian Institute of Information Technology Kalyani}, is an
authentic record of my own work carried out during the time period
from \textbf{July 2022} to \textbf{May 2025}. Thesis submission
under the supervision of ``\textbf{Name of the Supervisor,
Designation, and Affiliation}.

\noindent The matter presented in this thesis has not been submitted
by me for the award of any other degree of this or any other
institute.


\vspace{1in}

\begin{minipage}[l]{0.2\textwidth}%
\textbf{Date:}~
\\
\\
\\
\\
\\
\end{minipage}
\hfill
\begin{minipage}[r]{0.7\textwidth}%
\begin{flushright}
    \textbf{(Name of the student)}\\
    \textbf{Reg. No.:}~\\
    \text{Department of Computer Science and Engineering}\\
    \text{Indian Institute of Information Technology Kalyani}\\
    \text{West Bengal - 741235, India}
\end{flushright}
\end{minipage}

 \cleardoublepage
  \vspace*{0.1cm}
\begin{center}
{\LARGE \textsc{Acknowledgments}}
\end{center}
 \vspace{0.3cm}

%\begin{flushright}
%\begin{tabular}{l}
%"\textit{The time has come,}" \textit{the Walrus said,}\\
%"\textit{To talk of many things}" \dots \\
%\hspace{1in} - {\scriptsize Lewis Carroll, Through the Looking-Glass}\\
%\end{tabular}
%\end{flushright}

I am privileged to have the opportunity to carry out my..........\\

\vspace{1in}

\begin{minipage}[l]{0.2\textwidth}%
\textbf{Date:}~
\\
\\
\\
\\
\\
\end{minipage}
\hfill
\begin{minipage}[r]{0.7\textwidth}%
\begin{flushright}
    \textbf{(Name of the student)}\\
    \textbf{Reg. No.:}~\\
    \text{Department of Computer Science and Engineering}\\
    \text{Indian Institute of Information Technology Kalyani}\\
    \text{West Bengal - 741235, India}
\end{flushright}
\end{minipage}

 \cleardoublepage
  \cleardoublepage
 \pagenumbering{roman}
 \setcounter{page}{5}
 \begin{center}
 {\LARGE \textsc{Abstract}}
 \end{center}
 \vspace{1cm}
 Finish the Abstract with in one page\\
 \\
 \vspace{0.5cm} \noindent \textbf{Keywords:}~Provide 8-10 important keywords.

 %\cleardoublepage
 \tableofcontents
 \cleardoublepage
 \listoffigures
 \cleardoublepage
 \listoftables
 \cleardoublepage
 \listofalgorithms
 \cleardoublepage
 \newpage
%\begin{acronym}

\thispagestyle{plain}
%\phantomsection
%\label{acm}

\noindent\textbf{\huge List of Acronyms}
\vspace{0.2cm}
\setlength{\tabcolsep}{4pt}
\vspace{-1em}
\addtolength\tabcolsep{.01pt}
\renewcommand{\arraystretch}{1.325}
\setlength{\LTcapwidth}{12in}

%\begin{table}[!h]
   %\centering
   \begin{longtable}{p{5cm} p{9cm}}
   \hline
   \textbf{Abbreviation}  & \textbf{Full form}\\
   \hline
   \endfirsthead
    \hline
   \textbf{Abbreviation}  & \textbf{Full form}\\
   \hline
    \endhead

    \hline
    \endfoot
    \hline
    \endlastfoot
BDH     & Bilinear Diffie-Hellman (BDH) \\

\hline
\end{longtable}
%\end{table}

 \addcontentsline{toc}{chapter}{List of Acronyms}
 \clearpage

    \addtolength{\headheight}{15pt}
    \pagenumbering{arabic}
    %\setcounter{page}{1}
    \pagestyle{myheadings}
    %\markboth{Introduction}{Introduction}
    \pagestyle{fancy}
    \renewcommand{\chaptermark}[1]{\markboth{\chaptername
            \ \thechapter:\,\ #1}{}}
    \renewcommand{\sectionmark}[1]{\markright{\thesection\,\ #1}}
    \fancyhead{}
    %\fancyhead[LE]{\sl\leftmark}
    %\fancyhead[LO,RE]{\rm\thepage}
    %\fancyhead[RO]{\sl\rightmark}
    \fancyhead[LO]{\sl\rightmark}
    \fancyhead[LE,RO]{\rm\thepage}
    \fancyhead[RE]{\sl\leftmark}
    \fancyfoot[C,L,E]{}
%%----------------------GLOSSARY--------------------------------%
\clearpage
%%%%%%%%%%%%%%%%%%%%%%%%%%%%%%%%%%%%%%%%%%%%%%%%%%%%%%%%%%%%%%%%%%%%%%%%%%%%
%% Chapter 1
%%Indian Institute of Information Technology Kalyani
%% All rights are reserved.
%%%%%%%%%%%%%%%%%%%%%%%%%%%%%%%%%%%%%%%%%%%%%%%%%%%%%%%%%%%%%%%%%%%%%%%%%%%%
%

\chapter[Introduction]{Introduction}
\label{chp1}
%\setcounter{minitocdepth}{1}
%\minitoc
\section{Introduction}
\label{chp1.1}
 %%%%%%%%%%%%%%%%%%%%%%%%%%%%%%%%%%%%%%%%%%%%%%%%%%%%%%%%%%%%%%%%%%%%%%%%%%%%%%%%%%%%%%%
%%%%%%%%%%%%%%%%%%%%%%%%%%%%%%%%%%%%%%%%%%%%%%%%%%%%%%%%%%%%%%%%%%%%%%%%%%%%%%%%%%%%%%%
\begin{figure}[!htb]
\begin{center}  \fbox{\includegraphics[width=14.5cm,height=8.2cm]{Figure/chp1/chap1_1}}
\caption{General framework of cloud storage and secure data access}
\label{Fig1(Chp1)}
\end{center}
\end{figure}
%%%%%%%%%%%%%%%%%%%%%%%%%%%%%%%%%%%%%%%%%%%%%%%%%%%%%%%%%%%%%%%%%%%%%%%%%%%%%%%%%%%%%%%
%%%%%%%%%%%%%%%%%%%%%%%%%%%%%%%%%%%%%%%%%%%%%%%%%%%%%%%%%%%%%%%%%%%%%%%%%%%%%%%%%%%%%%%


\section{Background of PEKS Scheme}
\label{chp1.2}

\section{Research Scopes}

\section{Objectives of the Thesis}


\section{Contributions of the Thesis}

\section{Organization of the Thesis}

\section{Summary}

\cleardoublepage
\input{chp2/chp2}
\cleardoublepage
\cleardoublepage
%\begin{flushleft}
%{\Large \textsc{\textbf{Publications}}}
%\end{flushleft}
\chapter*{Publications}
%\rule{\textwidth}{0.4pt}\par % Thick horizontal line
This thesis is based on some of the following publications of the author.\\

{\noindent \Large \textbf{Journal Publications}}
\begin{enumerate}
\item xxxxxxx
\end{enumerate}

\addcontentsline{toc}{chapter}{Publications}
%\printglossary[style=long]
\bibliographystyle{ieeetr}
\cleardoublepage
\bibliography{Reference/Bibliography.bib}
\end{document}
